\section{Finite-free chain complexes}

The structure \code{IntChainComplex} stores a total rank function,
a finite-support bound, and matrices
\[
  \code{boundary}\ n :
  M_{\rank C_n,\rank C_{n+1}}(\Z).
\]
Thus \code{boundary n} always represents
$d_{n+1}\colon C_{n+1}\to C_n$ under the column-vector convention.  The
field \code{boundary\_squared} states that adjacent matrices multiply to
zero.

The total rank function makes every degree well typed.  Above
\code{topDegree} it is zero, and degree zero receives the unique outgoing
matrix of shape $0\times\rank C_0$.  Derived definitions distinguish:
\[
  \code{outgoingBoundary}\ n=d_n,\qquad
  \code{incomingBoundary}\ n=d_{n+1}.
\]
The theorem \code{outgoing\_mul\_incoming} covers both positive degrees and
the degree-zero empty matrix.

Using mathlib linear maps, we define
\[
\begin{aligned}
  Z_n&=\ker(d_n),\\
  B_n&=\operatorname{range}(d_{n+1}\colon C_{n+1}\to Z_n),\\
  H_n&=Z_n/B_n.
\end{aligned}
\]
The codomain restriction to $Z_n$ is justified directly by the matrix chain
condition.  Consequently the theorem-facing object is ordinary homology
before any coordinate calculation begins.

For small cellular examples,
\code{ofLengthTwo} constructs a complex concentrated in degrees zero through
two from matrices $d_1$, $d_2$, and a proof $d_1d_2=0$.  This constructor is
not a replacement for general bounded complexes; it makes the orientation
and regression fixtures concise.
